\documentclass[10pt,a4paper]{amsart}
\usepackage[margin=19mm]{geometry}
\usepackage{amsmath,amssymb,mathtools}
\usepackage[scr=rsfs,cal=euler]{mathalfa}
\usepackage{xfrac}
\usepackage[unicode,hidelinks,bookmarks=false]{hyperref}
\newcommand{\ie}{i.e.\ }
\newcommand{\cf}{cf.\ }
\newcommand{\resp}{respectively\ }
\newcommand{\Subsection}{\subsection}
\providecommand{\nobreakdash}{\nobreak\hbox{-}}
\setlength{\emergencystretch}{2em}
\multlinegap0pt
\usepackage{amssymb}
\usepackage{mathtools}
\usepackage{float}
\usepackage{enumerate}
\newtheorem{theorem}{Theorem}
\newtheorem{proposition}{Proposition}
\renewcommand{\L}{\mathrm{L}}
\newcommand{\Sl}{\mathrm{Sl}}
\newcommand{\So}{\mathrm{So}}
\newcommand{\U}{\mathrm{U}}
\newcommand{\cC}{\mathcal{C}}
\newcommand{\cO}{\mathcal{O}}
\title{Diagonalization of the Toda flow for \\ arbitrary isospectral symmetric matrices}
\author{David Mart\'inez Torres, Carlos Tomei}
\date{Source: arXiv:2608.07263v2 (CC BY 4.0)}
\begin{document}
\maketitle

\noindent\emph{Source and licence.} Diagonalization of the Toda flow for \\ arbitrary isospectral symmetric matrices. Version arXiv:2608.07263v2 (CC BY 4.0), distributed by arXiv under the Creative Commons Attribution 4.0 International licence, http://creativecommons.org/licenses/by/4.0/, as recorded in the arXiv licence metadata for this exact version and shown on the abstract page https://arxiv.org/abs/2608.07263v2. Full licence notice: \texttt{licenses/arxiv-2608.07263-CC-BY-4.0.txt}.\par\smallskip

\noindent\textit{Editorial scope.} Counted unit: the article's theorem thm:main (source0.tex lines 284--292). The article's complete original proof of it is delivered with it (source0.tex lines 554--589, one \texttt{proof} environment of 36 lines, which the article places away from the statement). No mathematics is written, repaired or re-derived: the statement and the proof are the article's own lines, in the article's own notation, and no retained line is substituted or re-flowed.  Retained uncounted as the source-local context the proof needs: lines 412--413; lines 419--425.  Nothing was omitted from any retained span, so the package carries no omission record.  The retained lines cite 1 works, whose entries are transcribed verbatim from the article's own bibliography source in the e-print and delivered with short editorial labels recorded in the item's reference mapping.  No retained line contains a figure environment, an image-inclusion command or drawing code, so no image is delivered and the package declares no asset dependency.  Editorial formatting only, in addition: an AMS \texttt{amsart} preamble that restates the article's own declarations and macros used by the retained lines, namely the environments \texttt{theorem}, \texttt{proposition} on separate counters, together with the standard packages those lines need.  The retained environments print in this document as Theorem 1, Proposition 1 in order of appearance, whereas the article prints the same items with the numbering of its own full document.  The complete native source of the article as distributed in the arXiv e-print for this exact version is bundled unchanged as \texttt{sources/207068/source0.tex}.
\begin{equation}\label{eq:diagonal-linear}L'=[L,-\Lambda],\quad L\in \L,
\end{equation}

\begin{proposition}\label{pro:Toda}
The map $\Psi: \L\to \cO $ takes  the trajectories of the action of $\mathrm{e}^{t\Lambda}$ on $\L$ by conjugation to the trajectories of the Toda vector field,
\[\Psi(\mathrm{e}^{t\Lambda}L\mathrm{e}^{-t\Lambda})=\kappa(\mathrm{e}^{tX})^T X  \kappa(\mathrm{e}^{tX}),
\quad X=\kappa(L)^T\Lambda \kappa(L).\]
Equivalently, the vector field  $L'=[L,-\Lambda]$  on $\L$
can be pushed forward by $\Psi$ and its image is the Toda vector field.
\end{proposition}

\begin{theorem}\label{thm:main} Let  $\cO$ be an orthogonal conjugacy class of real symmetric matrices.
	Then around each diagonal matrix $\Lambda\in \cO$ there exist local coordinates mapping onto $\L^\Lambda$
that transform the Toda vector field in the  diagonal linear  vector field
  \begin{equation}\label{eq:Toda-coordinates}
  L'=[L,-\Lambda],\quad L\in \L^\Lambda.
  \end{equation}

The domain of definition of each set of local coordinates is dense, and the union over the (finitely many) diagonal matrices in $\cO$ of these domains covers $\cO$.
\end{theorem}

\begin{proof}[Proof of Theorem \ref{thm:main}]
Suppose that the constrained $\mathrm{Q}\L\U$ factorization holds. The first projection of the $\mathrm{QR}$ factorization  $\kappa:\Sl\to  \So$ to $\So_\Lambda\L^\Lambda\U$ yields
 \[\kappa(\So_\Lambda\L^\Lambda\U)=\So_\Lambda \kappa(\L^\Lambda),\]
a factorization of $\So_\Lambda$ defined  in  the complement of the solution set of a polynomial equation (the restriction to $\So_\Lambda$ of homogeneous polynomial coming from the constrained $\mathrm{Q}\L\U$ factorizaton). Notice that   $\So_\Lambda \kappa(\L^\Lambda)$ is different from  $\cC=\kappa(\L)$.

As
\begin{itemize}\item $\kappa:\L^\Lambda\to\kappa(\L^\Lambda)$ is a diffeomorphism,   \item $\Psi$ is the result of applying anti-conjugation to $\Lambda$ by elements in $\kappa(\L^\Lambda)$,\item   $\So_\Lambda$ is the centralizer of $\Lambda$ in $\So$,
 \item
 the (unique and smooth) factorization $\So_\Lambda \kappa(\L^\Lambda)\subset \So$ holds,
 \end{itemize}
 the same argument that we sketched for the simple spectrum case implies
that
\[\Psi:\L^\Lambda\longrightarrow\cO\]
is a diffeomorphism onto its image $\mathcal{V}\subset \cO$. Because $\So_\Lambda\kappa(\L^\Lambda)\subset \So$ is open and dense, then so is $\mathcal{V}\subset \cO$.

As $\L^\Lambda$ is a coordinate affine subspace of $\L$, the diagonal linear vector field  \eqref{eq:diagonal-linear} is tangent to
        $\L^\Lambda$. Therefore
Proposition \ref{pro:Toda} implies that in the local coordinates $\Psi^{-1}:\mathcal{V}\to \L^\Lambda$  the Toda vector field equals \eqref{eq:diagonal-linear}.

%Next, we discuss additional properties of the domains of the local coordinates.
%The  vector field \eqref{eq:diagonal-linear} on $\L^\Lambda$ has a unique zero at the origin, since
%\[L_{ij}'=(\lambda_i-\lambda_j)L_{ij},\quad \lambda_i-\lambda_j\neq 0.\]
%This means that the only diagonal matrix in $\mathcal{V}$ is $\Lambda$.
% Also, because \eqref{eq:diagonal-linear} is  complete $\mathcal{V}$ contains both the stable and unstable manifold at $\Lambda$ of the Toda vector field.

We argue that
$\cO$ is also covered by the union of the domains of the local coordinates centered at each diagonal matrix by using a result from dynamical systems: If a vector field on compact manifold has hyperbolic zeroes and has a strict Lyapunov function, then the union of its stable (or unstable) manifolds over its zeroes covers the manifold. Back to our setting,
\begin{itemize}
 \item the conjugacy class $\cO$ is compact,
 \item the Toda vector field on $\cO$ has hyperbolic zeroes, the diagonal matrices in the conjugacy class,
 \item  in the linear coordinates centered at a diagonal matrix  the stable manifold is a vector subspace and therefore it sits in the domain of the coordinates,
 \item the restriction to $\cO$ of the linear function
 \[f(X)=\mathrm{tr}(AX)\]
 where a is a fixed diagonal matrix with different strictly positive entries ordered increasingly, is a strict Lyapunov function for the Toda vector field  \cite[Lemma 1]{Fa}.
\end{itemize}
\end{proof}

\begin{thebibliography}{99}
\bibitem[Fa]{Fa} Faybusovich, L. Toda flows and isospectral manifolds. Proc. Amer. Math. Soc. 115 (1992), no. 3, 837--847.
\end{thebibliography}
\end{document}
